\documentclass[xcolor=dvipsnames, compress, 12pt, aspectratio=169, handout]{beamer}
\usepackage{../preamble_slides}
\usepackage{pifont}

\title{Labor Market Effects of Winning the H-1B Lottery}
\author{Amy Kim}
\date{September 11, 2025}

\begin{document}
\begin{frame}[plain]
    \titlepage
    \addtocounter{framenumber}{-1}
\end{frame}

\begin{frame}
    \frametitle{Project Overview}
    \begin{itemize}
        \item \textbf{Research Question:} What is the individual-level effect of winning the H-1B lottery on labor market outcomes?
        \vspace{2mm}
        \item \textbf{Data:} H-1B lottery applications and LinkedIn profile data from Revelio
        \vspace{2mm}
        \item \textbf{Empirical challenge:} linkage with limited information 
        \begin{itemize}
            \item H-1B data contains only country of nationality, year of birth, gender, firm 
            \item But hard to use firm\textemdash if individual loses lottery, can't observe as working at firm?
        \end{itemize}
        \vspace{2mm}
        \item \textbf{Solution:} Restrict to sample of international STEM students 
        \begin{itemize}
            \item STEM grads have 3-year work permit (OPT)
            \item Hypothesis: STEM grads work at firm on OPT, apply for H-1B lottery every year until they win
            \item Should therefore be able to observe as working at firm regardless of lottery outcome
        \end{itemize}
    \end{itemize}
\end{frame}

\begin{frame}
    \frametitle{Project Progress}
    \begin{itemize}
        \item This summer: 
        \begin{itemize}
            \item First pass at linkage, much harder than anticipated 
            \item Initial results difficult to interpret (real or just noise?)
            \item \textbf{Conclusion:} this paper will be a noisy/shaky stab at a very big and important question (not ideal for JMP but still meaningful contribution to field)
            \item Next steps: write up paper ASAP and pivot
        \end{itemize}
        \item \textbf{Today:} present results, talk about obstacles
        \item Looking for feedback on:
        \begin{itemize}
            \item Most compelling `story' for this paper (with minimal additional work)
            \item Empirics: inference with probabilistic matching 
            \item Interesting future directions/projects 
        \end{itemize}
    \end{itemize}
\end{frame}

\begin{frame}
    \frametitle{Step One: Linkage}
    \textbf{Overview}
    \begin{itemize}
        \item I impute nationality, YOB, gender for individual LinkedIn users and match to H-1B applications, assigning weights based on how certain I am of a match
        \item Then I further filter matches based on tenure at firm, occupation history, etc.
        \item \textbf{Resulting Baseline Sample:} 989k LinkedIn users matched to 78k applications (`multiplicity' of 13)
        \begin{itemize}
            \item Alt. Samples: restrict the sample to exclude common countries, larger firms, etc.
        \end{itemize}
    \end{itemize}
    \vspace{2mm}
    \textbf{Key Empirical Problems}
    \begin{enumerate}
        \item Selection/External Validity: lower match weight for individuals with less info on profile, no benchmark for how many applications I \textit{should} be able to match
        \item Noise: Each application matched to many potential users, dilutes true estimate
    \end{enumerate}
\end{frame}

\begin{frame}
    \frametitle{Step Two: Collapsing Matches}
    \textbf{Overview}
    \begin{itemize}
        \item For now: I collapse observations to the application level (compute weighted average of relevant variables using whatever weights I have)
        \item Eventually: calibrate weights + simulate draws based on weights (`bootstrap'-style) 
    \end{itemize}
    \vspace{2mm}
    \textbf{Key Empirical Problem:}
    \begin{enumerate}
        \item Simulating draws is complicated because many LinkedIn users are matched (with high weight) to multiple potential applications\textemdash prior literature? Inference?
    \end{enumerate}
\end{frame}
\begin{frame}
    \frametitle{Step Three: Estimation}
    For application $j$ for an individual at company $c$ in lottery year $t$ I estimate:
    \begin{equation}
        y_{jct} = \beta_0 + \beta_1 win_j + \beta_2 ADE_j + \varepsilon_{jct}
    \end{equation}
    where $win_j$ indicates winning the lottery and $ADE_j$ indicates `Advanced Degree Exemption' status (the only thing that changes your lottery odds), outcomes $y_{jct}$ include:
    \begin{itemize}
        \item $work1yr$ ($work2yr$): indicator for still working at firm one (two) years after lottery
        \item $us1yr$ ($us2yr$): indicator for still working in the U.S. one (two) years after lottery
        \item $educ1yr$ ($educ2yr$): indicator for re-entering school one (two) years after lottery
    \end{itemize}
    In some specifications I also include firm x year fixed effects.

    \vspace{2mm}
    \textbf{Identification Assumption:} winning is random conditional on ADE status \textemdash this seems to be true! Yay!
\end{frame}

\begin{frame}
    \frametitle{Results (one year post lottery)}
    \centering 
    \resizebox{\textwidth}{!}{
\begin{tabular}{lcccccc}
\toprule
& \multicolumn{2}{c}{work1yr} & \multicolumn{2}{c}{us1yr} & \multicolumn{2}{c}{educ1yr} \\
& (1) & (2) & (3) & (4) & (5) & (6) \\
\midrule
win & 0.0186*** & 0.0054** & 0.0071*** & 0.0021* & -0.0082*** & 0.0006 \\
& (0.0024) & (0.0025) & (0.0011) & (0.0012) & (0.0024) & (0.0026) \\
\midrule
Firm x Year FEs & No & Yes & No & Yes & No & Yes \\
\midrule 
\textbf{Obs.} & 77,606 & 77,606 & 77,606 & 77,606 & 77,606 & 77,606 \\
\textbf{DV mean} & 0.690 & 0.690 & 0.950 & 0.950 & 0.248 & 0.248 \\
\bottomrule
\end{tabular}
    }
\end{frame}

\begin{frame}
    \frametitle{Results (two years post lottery)}
    \centering 
    \resizebox{\textwidth}{!}{
\begin{tabular}{lcccccc}
\toprule
& \multicolumn{2}{c}{work2yr} & \multicolumn{2}{c}{us2yr} & \multicolumn{2}{c}{educ2yr} \\
& (1) & (2) & (3) & (4) & (5) & (6) \\
\midrule
win &  0.0052* & 0.0010 & 0.0153*** & 0.0038** & -0.0049** & 0.0029 \\
& (0.0027) & (0.0027) & (0.0015) & (0.0016) & (0.0025) & (0.0027) \\
\midrule
Firm x Year FEs & No & Yes & No & Yes & No & Yes \\
\midrule 
\textbf{Obs.} & 77,606 & 77,606 & 77,606 & 77,606 & 77,606 & 77,606 \\
\textbf{DV mean} & 0.523 & 0.523 & 0.911 & 0.911 & 0.276 & 0.276 \\
\bottomrule
\end{tabular}
    }
\end{frame}

\begin{frame}
    \frametitle{Benchmarking Results: Why Are Coefficients So Small?}
    When estimating the coefficient $\beta_1$ on $win$ for outcome $work1yr$, we can decompose:
    \begin{equation*}
        \mathbb{E}[\hat{\beta}_1] = r \cdot s_3 \cdot w \approx 0.07
    \end{equation*}
    \begin{itemize}
        \item $r \equiv \mathbb{E}[work1yr_{jct}|win_j = 1]$ represents the probability of an individual staying at the firm conditional on winning the lottery 
        \begin{itemize}
            \item Estimate in data: $r \approx 0.7$
        \end{itemize}
        \item $s_{last}$ represents the share of applicants that are on their `last attempt' at the lottery
        \begin{itemize}
            \item Estimate in theory: if all applicants get three tries and win probability is 0.4 in all years, $s_{last} \approx 0.18$
        \end{itemize}
        \item $w$ represents the average weight of the `true' applicant among all potential matches
        \begin{itemize}
            \item Estimate in data (mean of max weight): $w \approx 0.54$ 
        \end{itemize}
    \end{itemize}
\end{frame}

\begin{frame}
    \frametitle{Conclusion and Next Steps}
    \begin{itemize}
        \item \textbf{Fundamental Problem:} difficult to distinguish whether coefficient small due to $s_{last}$ (`real' policy-relevant parameter) or $w$ (measurement error)
        \vspace{2mm}
        \item What outcomes/channels are most compelling to dig further into? 
        \begin{itemize}
            \item Promotions (change jobs within firm)
            \item Entrepreneurship/mobility (changing firms, focus on smaller firms)
            \item Education (idk but seems weird anecdotally that there's not much of a result)
            \item Conditional on leaving US: where do they go? what kinds of jobs do they work?
            \item If they lose lottery and don't leave firm, what does their trajectory look like?
        \end{itemize}
    \end{itemize}
\end{frame}

\appendix
\begin{frame}
    \frametitle{Empirical Strategy: Basic Version}
For H-1B lottery application $j$:
\begin{equation*}
    stay^{1yr}_{i(j), c(j), t(j)} = \beta_0 + \beta_1 win_{j} + \varepsilon{j}
\end{equation*}

where $i(j)$, $c(j)$, and $t(j)$ denote the individual, company, and lottery year respectively of application $j$, $stay_{i,c,t}$ indicates whether individual $i$ who was observed working at company $c$ in year $t$ is still working at the firm one year later (in $t+1$). $win_{j}$ indicates whether application $j$ was selected in the H-1B lottery.

\vspace{2mm}
Our coefficient of interest is $\beta_1$, the effect of winning the lottery on staying at the firm (suppressing some notation for convenience):
\begin{equation*}
    \beta_1 = \mathbb{E}[stay_{ict}|win_j = 1] - \mathbb{E}[stay_{ict}|win_j = 0]
\end{equation*} 
\vspace{2mm}

In a world where all lottery losers leave immediately ($\mathbb{E}[stay_{ict}|win_j = 0] = 0$) and everyone else stays with probability $r$ ($\mathbb{E}[stay_{ict}|win_j = 1] = r$), we have $\beta_1 = r$. 

\end{frame}

\begin{frame}
\frametitle{Complication 1: Losers Don't Leave Immediately}
In my setting, the first reason why I might be estimating a smaller $\beta_1$ than expected is because of the repeated draw structure of the lottery.
\vspace{2mm}

Let's assume that all applicants $i(j)$ are on 3-year OPT visas, so they apply for the H-1B visa every year until they win, with a max of three attempts. Applicants who are on their third try and lose leave the company immediately, otherwise they stay with probability $r$. Then:
\begin{align*}
    \mathbb{E}[stay_{ict}|win_j = 0] &= \mathbb{E}[\mathbb{E}[stay_{ict}|win_j=0, attempt(j) = a] | win_j = 0] \\
    &= \mathbb{P}[attempt(j) < 3] \cdot r + \mathbb{P}[attempt(j) = 3] \cdot 0 \\
    &= (1-s_3) \cdot r
\end{align*} 
where $attempt(j) \in \{1,2,3\}$ denotes the `attempt number' of application $j$ and $s_3$ represents the share of applicants that are on their third attempt.
\end{frame}

\begin{frame}
    \frametitle{Complication 2: Multiplicity/Mismeasurement}
    The second complication is when I measure $stay_{ict}$, I'm really measuring a weighted average of the true $stay_{ict}$ and the $stay_{i'ct}$ for some falsely matched individuals $i'$. If the true individual has weight $w$ on average and the falsely matched individuals have average stay rate $q$:

    \begin{align*}
        \mathbb{E}[\hat{\beta}_1] &= w\cdot r + (1-w) \cdot q - (w \cdot (1-s_3) \cdot r + (1-w) \cdot q) \\ 
        &= w s_3 r
    \end{align*}

    Note that all estimated coefficients (not just for outcome stay rate) will be downweighted by $w$ under these assumptions.
\end{frame}

\begin{frame}
    \frametitle{Putting it All Together}
    \begin{equation*}
        \mathbb{E}[\hat{\beta}_1] = w s_3 r
    \end{equation*}
    \begin{itemize}
        \item Estimating $r$ (average stay rate of winners): various methods yield similar estimates of 0.5
        \item Estimating $s_3$ (share of applicants in sample that are on their third try): assuming constant win probability of $0.4$ in sample yields estimate of 0.18
        \item Estimating $w$ (average weight on `true match'): depending on merge specification, estimates range from 0.5 to 0.7
    \end{itemize}
    \begin{align*}
        \hat{w} \cdot \hat{s}_3 \cdot \hat{r} \in [0.5 \cdot 0.18 \cdot 0.5, 0.7 \cdot 0.18 \cdot 0.5] = [0.045, 0.063]
    \end{align*}
    So even if I'm estimating and measuring everything correctly, in the best matched sample I would only expect to see a coefficient of $0.06$.
\end{frame}

\begin{frame}
    \frametitle{Some Implications}
    \begin{equation}
        \mathbb{E}[\hat{\beta}_1] = w s_3 r
    \end{equation}
    \begin{enumerate}
        \item The estimated coefficient on $stay_{ict}$ should increase as $s_3$ (the share of the sample that is on their third try) increases.
        \begin{itemize}
            \item I can vary this either by restricting the lottery year ($s_3$ should increase in later years as the lottery gets harder), or by measuring outcome $stay^{2yr}$ (the resulting coefficient would be downweighted instead by $s_2 > s_3$)
            \item If the estimated coefficient does not increase with $s_3$, it suggests that my ex-ante three year OPT model of behavior is incorrect (or that there is mismeasurement)
        \end{itemize}
        \item The estimated coefficient on $stay_{ict}$ should increase as $w$ (the average weight of the `true' individual) increases.
        \begin{itemize}
            \item I can vary this in the `postfiltered' merge sample, which is explicitly filtered on multiplicity
            \item If the estimated coefficient does not increase with $w$, it suggests mismeasurement
        \end{itemize}
    \end{enumerate}
\end{frame}

\begin{frame}
    \frametitle{Updates: Remain at Company After 1 Year}
    
    \begin{itemize} 
        \item Issue: $\beta_1$ `diluted' because of repeat tries at lottery 
        \item Solution: subset to sample less likely to try again
        \begin{itemize}
            \item First strategy: employers that exhibit `repeat try' behavior vs. employers that do not 
            \begin{itemize}
                \item Would expect `repeat try' employers sample to have much smaller coefficient than `non-repeat try' employers 
            \end{itemize}
            \item Second strategy: filter by years after graduation 
            \begin{itemize}
                \item Would expect lottery to matter more for people >3 years out of school
            \end{itemize}
        \end{itemize}
        \item Sample: Full Baseline sample (with high multiplicity but large sample size)
    \end{itemize}
\end{frame}


\begin{frame}
    \frametitle{Raw Data: Mean of stay1yr by Years Since Graduation}
    \centering
    \includegraphics[width=0.7\textwidth]{\outputpath/graddiff_work1yr_sep9.png}
\end{frame}

\end{document}
